% GENERATED FILE. Edit canonical lessons, then run npm run generate:books.
% canonical-source-hash: fnv1a64:58a6178d3c938e50
% canonical-lessons: KA-C211-niyamitavagi, KA-C211-illiyavarege, KA-C211-varege, KA-C211-sadyakke, KA-C211-avadhi

\chapter{Regularly, So far, Until, For now, Period}
\label{ch:ka-regularly-so-far-until-for-now-period}
\hlchaptermodality{211}

\begin{chapteropening}
\textbf{By the end of this chapter:} \emph{I can say regularly, so far, until, for now and a period.}

Complete the last lesson of chapter 211: I can say regularly, so far, until, for now and a period.
\end{chapteropening}

\section[\texorpdfstring{\kn{ನಿಯಮಿತವಾಗಿ} (niyamitavāgi)}{ನಿಯಮಿತವಾಗಿ (niyamitavāgi)}]{\kn{ನಿಯಮಿತವಾಗಿ} (niyamitavāgi) — regularly}

\label{lesson:KA-C211-niyamitavagi}



\begin{quote}
Before the new one: say the Kannada for free time, then the Kannada for Ugadi.
\end{quote}

\subsection*{You'll want to know: \kn{ನಿಯಮಿತವಾಗಿ}}
\textbf{\kn{ನಿಯಮಿತವಾಗಿ}} — \emph{niyamitavāgi} — \textquotedblleft{}regularly\textquotedblright{}.

\begin{grammarlens}[title={What you've built}]
One more word for this chapter.
\end{grammarlens}

\subsection*{Your turn}
\begin{itemize}
\raggedright
  \item \emph{Say it:} \emph{niyamitavāgi}
  \item \emph{Say it:} \emph{niyamitavāgi}, once more
  \item \emph{Say it:} say \emph{niyamitavāgi}
  \item \emph{Recall:} say the Kannada for free time, then the Kannada for Ugadi, then say \emph{niyamitavāgi} again
\end{itemize}

\begin{tcolorbox}[breakable,colback=teal!4,colframe=teal!35!black,title={Before you move on}]
What does \kn{ನಿಯಮಿತವಾಗಿ} mean? (\textquotedblleft{}Regularly\textquotedblright{}.) Say it once more.
\end{tcolorbox}
\section[\texorpdfstring{\kn{ಇಲ್ಲಿಯವರೆಗೆ} (illiyavarege)}{ಇಲ್ಲಿಯವರೆಗೆ (illiyavarege)}]{\kn{ಇಲ್ಲಿಯವರೆಗೆ} (illiyavarege) — so far, until now}

\label{lesson:KA-C211-illiyavarege}



\begin{quote}
Before the new one: say the Kannada for Ugadi, then the Kannada for regularly.
\end{quote}

\subsection*{You'll want to know: \kn{ಇಲ್ಲಿಯವರೆಗೆ}}
\textbf{\kn{ಇಲ್ಲಿಯವರೆಗೆ}} — \emph{illiyavarege} — \textquotedblleft{}so far, until now\textquotedblright{}.

\begin{grammarlens}[title={What you've built}]
One more word for this chapter.
\end{grammarlens}

\subsection*{Your turn}
\begin{itemize}
\raggedright
  \item \emph{Say it:} \emph{illiyavarege}
  \item \emph{Say it:} \emph{illiyavarege}, once more
  \item \emph{Say it:} say \emph{illiyavarege}
  \item \emph{Recall:} say the Kannada for Ugadi, then the Kannada for regularly, then say \emph{illiyavarege} again
\end{itemize}

\begin{tcolorbox}[breakable,colback=teal!4,colframe=teal!35!black,title={Before you move on}]
What does \kn{ಇಲ್ಲಿಯವರೆಗೆ} mean? (\textquotedblleft{}So far, until now\textquotedblright{}.) Say it once more.
\end{tcolorbox}
\section[\texorpdfstring{\kn{ವರೆಗೆ} (varege)}{ವರೆಗೆ (varege)}]{\kn{ವರೆಗೆ} (varege) — until, up to}

\label{lesson:KA-C211-varege}



\begin{quote}
Before the new one: say the Kannada for regularly, then the Kannada for so far.
\end{quote}

\subsection*{You'll want to know: \kn{ವರೆಗೆ}}
\textbf{\kn{ವರೆಗೆ}} — \emph{varege} — \textquotedblleft{}until, up to\textquotedblright{}.

\begin{grammarlens}[title={What you've built}]
One more word for this chapter.
\end{grammarlens}

\subsection*{Your turn}
\begin{itemize}
\raggedright
  \item \emph{Say it:} \emph{varege}
  \item \emph{Say it:} \emph{varege}, once more
  \item \emph{Say it:} say \emph{varege}
  \item \emph{Recall:} say the Kannada for regularly, then the Kannada for so far, then say \emph{varege} again
\end{itemize}

\begin{tcolorbox}[breakable,colback=teal!4,colframe=teal!35!black,title={Before you move on}]
What does \kn{ವರೆಗೆ} mean? (\textquotedblleft{}Until, up to\textquotedblright{}.) Say it once more.
\end{tcolorbox}
\section[\texorpdfstring{\kn{ಸದ್ಯಕ್ಕೆ} (sadyakke)}{ಸದ್ಯಕ್ಕೆ (sadyakke)}]{\kn{ಸದ್ಯಕ್ಕೆ} (sadyakke) — for now, for the time being}

\label{lesson:KA-C211-sadyakke}



\begin{quote}
Before the new one: say the Kannada for so far, then the Kannada for until.
\end{quote}

\subsection*{You'll want to know: \kn{ಸದ್ಯಕ್ಕೆ}}
\textbf{\kn{ಸದ್ಯಕ್ಕೆ}} — \emph{sadyakke} — \textquotedblleft{}for now, for the time being\textquotedblright{}.

\begin{grammarlens}[title={What you've built}]
One more word for this chapter.
\end{grammarlens}

\subsection*{Your turn}
\begin{itemize}
\raggedright
  \item \emph{Say it:} \emph{sadyakke}
  \item \emph{Say it:} \emph{sadyakke}, once more
  \item \emph{Say it:} say \emph{sadyakke}
  \item \emph{Recall:} say the Kannada for so far, then the Kannada for until, then say \emph{sadyakke} again
\end{itemize}

\begin{tcolorbox}[breakable,colback=teal!4,colframe=teal!35!black,title={Before you move on}]
What does \kn{ಸದ್ಯಕ್ಕೆ} mean? (\textquotedblleft{}For now, for the time being\textquotedblright{}.) Say it once more.
\end{tcolorbox}
\section[\texorpdfstring{\kn{ಅವಧಿ} (avadhi)}{ಅವಧಿ (avadhi)}]{\kn{ಅವಧಿ} (avadhi) — a period, a duration}

\label{lesson:KA-C211-avadhi}



\begin{quote}
Before the new one: say the Kannada for until, then the Kannada for for now.
\end{quote}

\subsection*{You'll want to know: \kn{ಅವಧಿ}}
\textbf{\kn{ಅವಧಿ}} — \emph{avadhi} — \textquotedblleft{}a period, a duration\textquotedblright{}.

\begin{grammarlens}[title={What you've built}]
One more word for this chapter.
\end{grammarlens}

\subsection*{Your turn}
\begin{itemize}
\raggedright
  \item \emph{Say it:} \emph{avadhi}
  \item \emph{Say it:} \emph{avadhi}, once more
  \item \emph{Say it:} say \emph{avadhi}
  \item \emph{Recall:} say the Kannada for until, then the Kannada for for now, then say \emph{avadhi} again
\end{itemize}

\begin{tcolorbox}[breakable,colback=teal!4,colframe=teal!35!black,title={Before you move on}]
What does \kn{ಅವಧಿ} mean? (\textquotedblleft{}A period, a duration\textquotedblright{}.) Say it once more.
\end{tcolorbox}
